\documentclass[12pt,a4paper]{amsart}

%% ── Packages ─────────────────────────────────────────────────────────────────
\usepackage{amsmath,amssymb,amsthm}
\usepackage{mathtools}
\usepackage[hidelinks]{hyperref}
\usepackage{cleveref}
\usepackage{enumitem}
\usepackage{booktabs}
\usepackage[margin=2.5cm]{geometry}

\emergencystretch=3em

\crefname{problem}{Open Problem}{Open Problems}
\Crefname{problem}{Open Problem}{Open Problems}

%% ── Theorem environments ──────────────────────────────────────────────────────
\theoremstyle{plain}
\newtheorem{theorem}{Theorem}[section]
\newtheorem{lemma}[theorem]{Lemma}
\newtheorem{corollary}[theorem]{Corollary}
\newtheorem{proposition}[theorem]{Proposition}

\theoremstyle{definition}
\newtheorem{definition}[theorem]{Definition}
\newtheorem{remark}[theorem]{Remark}
\newtheorem{problem}[theorem]{Open Problem}

%% ── Macros ────────────────────────────────────────────────────────────────────
\newcommand{\R}{\mathbb{R}}
\newcommand{\Tr}{\operatorname{Tr}}

% Canonical system Hamiltonian
\newcommand{\Ham}{H}
% de Branges structure function
\newcommand{\dbE}{E}

%% ── Scope warning macro ───────────────────────────────────────────────────────
\newcommand{\scopenote}[1]{\par\smallskip\noindent\textbf{Scope note.} #1\par\smallskip}

%% ─────────────────────────────────────────────────────────────────────────────

\title[A Limit-Point Criterion and Obstacles to a dBN Realization]{A Limit-Point Criterion for Canonical Systems\\
       and Obstacles to a de Bruijn--Newman Realization}

\author{Lin Tao}
\address{Shanghai}
\date{\today}

\begin{document}

\begin{abstract}
We record, with a short proof for reference, the standard trace-divergence
criterion (\cref{prop:weyl}) for the limit-point classification of the
right endpoint of a two-dimensional canonical system, under Remling's
standing assumptions; it is a direct specialisation of Remling's
Theorem~3.5 and is not a new result.
We then audit the ``Core Lemma'' of the speculative programme
\cite{SpecRoadmap}.  Two findings are negative.  First, the naive
realization via the logarithmic-derivative Herglotz function
$m_t=-H_t'/H_t$ exists but is too weak to be useful: it produces a
trace-normalised half-line system that is automatically limit-point,
making the endpoint conclusion tautological (\cref{sec:trivial}).
Second, and more importantly, the programme's proposed equivalence
``$\mathrm{RH}\iff\operatorname{spec}T_\Phi=\{\gamma_n\}$'' is \emph{not
justified as written}: a self-adjoint operator has real spectrum, so its
eigenvalues $\{\gamma_n\}$ (the ordinates of the zeta zeros) record
nothing about the \emph{real parts} of those zeros, and the roadmap
supplies no theorem recovering that missing information (\cref{sec:open}).
Since $T_\Phi$ has not been rigorously defined, we do not claim to have
disproved the implication for that specific relation; we show only that
the stated justification is insufficient.  A correct sufficient
replacement is a one-directional criterion using the characteristic
function $\Xi(z)=\xi(\tfrac12+iz)$; it is at least as hard as RH, imposing
requirements not known to follow from RH.  The paper yields \emph{no}
nontrivial or arithmetically meaningful theorem about the de Bruijn--Newman
flow.

\scopenote{This paper is a reference note plus a logical-gap and
proof-obligation audit of a speculative programme.  It does \emph{not}
imply the Riemann Hypothesis or any consequence thereof.  The only proved
statements are the criterion \cref{prop:weyl} and the trivial realization
of \cref{sec:trivial}; the programme's Core Lemma is shown to be, with
only the justification it gives, insufficient for RH, and the corrected
problem \cref{prob:gen} is marked \textbf{OPEN}.}
\end{abstract}

\maketitle

\tableofcontents

%% ─────────────────────────────────────────────────────────────────────────────
\section{Introduction}
\label{sec:intro}

We work throughout in the Rodgers--Tao convention \cite{RodgersTao2020}.
Let
\[
  \Phi_{\mathrm{RT}}(u) \;=\; \sum_{n\ge 1}
    \bigl(2\pi^2 n^4 e^{9u} - 3\pi n^2 e^{5u}\bigr)\,e^{-\pi n^2 e^{4u}},
    \qquad u \in \R,
\]
and define the heat-flow deformation, for $t \in \R$, by
\begin{equation}\label{eq:heatflow}
  H_t(z) \;=\; \int_0^\infty e^{t u^2}\,\Phi_{\mathrm{RT}}(u)\cos(zu)\,du,
  \qquad
  \partial_t H_t(z) \;=\; -\,\partial_z^2 H_t(z),
\end{equation}
so that $H_0(z) = \tfrac18\,\xi\!\bigl(\tfrac12 + \tfrac{iz}{2}\bigr)$.
Note the \emph{backward} heat equation carries a minus sign; it is the
kernel weight $e^{tu^2}$ (not a Gaussian multiplier in the transform
variable) that genuinely moves the non-real zeros of $H_t$ toward the
real axis as $t$ increases.  In particular the naive multiplier
$e^{-t z^2} H_0(z)$ would leave the zero set of $H_0$ unchanged and does
\emph{not} realise the de Bruijn--Newman flow.
The de Bruijn--Newman constant $\Lambda \in [0, 0.2]$ (lower bound
$\Lambda \ge 0$ by Rodgers--Tao \cite{RodgersTao2020}; upper bound
$\Lambda \le 0.2$ by Saouter--Gourdon--Demichel \cite{SaouterGourdon2011},
using zero-free region data from Platt--Trudgian \cite{PlattTrudgian2021})
is defined by:
for $t > \Lambda$, $H_t$ has only real zeros
(Laguerre--P\'olya class; cf.\ Newman \cite{Newman1976}); for
$t < \Lambda$, $H_t$ has non-real zeros.
The Riemann Hypothesis is equivalent to $\Lambda = 0$.

\begin{remark}[Relation to the earlier scaling]
  A previous draft used the Titchmarsh-scaled kernel
  $\Phi_{\mathrm{doc}}(x) = 2\,\Phi_{\mathrm{RT}}(x/2)$ together with a
  Gaussian factor $e^{-\lambda t^2}\Xi_0(t)$.  That factor is nowhere
  zero and hence leaves the zero set of $\Xi_0$ unchanged, so it does not
  implement the flow.  Under the correct kernel identity one has
  $\int_0^\infty e^{\lambda x^2}\Phi_{\mathrm{doc}}(x)\cos(tx)\,dx
   = 4\,H_{4\lambda}(2t)$, whence $\Xi_0(t) = \tfrac12\,\xi(\tfrac12+it)$
  and the critical parameter rescales to $\Lambda/4$.  To avoid these
  pitfalls we use only the Rodgers--Tao normalisation above.
\end{remark}

De Branges' theory \cite{deBranges1968} associates to a $2\times 2$
positive-semidefinite locally-integrable Hamiltonian $\Ham(x)$ the
canonical system
\begin{equation}\label{eq:cansys}
  J\,y'(x) = z\,\Ham(x)\,y(x), \qquad
  J = \begin{pmatrix} 0 & -1 \\ 1 & 0 \end{pmatrix}.
\end{equation}
At each regular truncation $x$ the transfer matrix yields a de~Branges
structure function $E_x$ and space $\mathcal H(E_x)$, whose reproducing
kernel is an \emph{internal} kernel of that space.  After a self-adjoint
realization is fixed, the associated spectral transform maps the operator
(or relation) part unitarily onto an $L^2(\R,\rho)$ space for a spectral
measure $\rho$; the Krein--de~Branges correspondence
\cite{Krein1952,Remling2018} relates $\Ham$, the chain $\{E_x\}$, and
$\rho$.  We stress that $\rho$ is \emph{not} the reproducing kernel:
the two are distinct objects.

The key structural question for the programme of \cite{SpecRoadmap} is:
\emph{does a canonical system attached to the de Bruijn--Newman flow
satisfy the limit-point condition at $x=+\infty$?}  A limit-point right
endpoint removes the boundary freedom at that endpoint; it does \emph{not}
by itself make the associated relation an essentially self-adjoint
operator, since the regular left endpoint $x=0$ still carries one boundary
condition (see \cref{prop:weyl} and \cref{cor:alpha}).

\subsection*{What this paper does and does not do}

The D1 step of \cite{SpecRoadmap} gave a \emph{qualitative} argument.
The present paper isolates what is actually provable, and is best read as
a research audit rather than a source of new theorems.

\begin{enumerate}[label=(\arabic*)]
  \item \textbf{Reference criterion (\cref{sec:L3}).}  We restate, with a
        short proof, the standard trace-divergence limit-point criterion
        \cref{prop:weyl}.  This is a direct specialisation of
        \cite[Theorem~3.5]{Remling2018}; it is not new and its proof
        depends on \cite{Remling2018}.
  \item \textbf{Trivial realization (\cref{sec:trivial}).}  For
        $t>\Lambda$, the logarithmic derivative $m_t=-H_t'/H_t$ is a
        generalised Herglotz function, so Remling's inverse theorem
        \cite[Theorem~5.1]{Remling2018} produces a \emph{unique
        trace-normalised} half-line canonical system $\widehat\Ham_t$.
        Being trace-normalised, it has
        $\int_0^\infty \Tr\widehat\Ham_t = +\infty$ automatically and is
        therefore limit-point by \cref{prop:weyl} --- a tautology that
        carries no arithmetic content.
  \item \textbf{Audit of the Core Lemma (\cref{sec:open}).}  According to
        \cite{SpecRoadmap} (its ``Core Lemma~L''), D1's acceptance
        criterion is a two-part equivalence: $T_\Phi$ limit-point at
        $+\infty$ \emph{and} a fixed self-adjoint realization having
        spectrum $\{\gamma_n\}$.  We find its stated justification
        \emph{insufficient as written}: the second conjunct records only
        the ordinates $\gamma_n$, which are real regardless of RH, so
        self-adjointness plus equality of ordinate sets does \emph{not} by
        itself imply RH; the roadmap gives no theorem recovering the
        real-part information (\cref{rem:coreL-defective}).  We do
        \emph{not} disprove the implication for the (still undefined)
        $T_\Phi$.  A correct sufficient replacement uses the
        characteristic function $\Xi(z)=\xi(\tfrac12+iz)$: a full
        zero-divisor identification with $\Xi$ \emph{does} imply RH, as a
        one-way sufficiency at least as hard as RH
        (\cref{prob:gen,rem:coreL-is-RH}).  We also note that $T_\Phi$ and
        its Hamiltonian are not yet rigorously defined in
        \cite{SpecRoadmap}; constructing them is the first, unstated part
        of the problem (\cref{rem:undefined}).  No terminal entire
        function at $x=\infty$ is assumed (\cref{rem:incompat}).
\end{enumerate}

Consequently the paper contains \emph{no} arithmetically meaningful
limit-point theorem about the de Bruijn--Newman flow, and nothing that
bears on RH.

\begin{remark}[Trace under reparametrisation]\label{rem:invariance}
  Under the admissible reparametrisation $X(x)=\int_0^x \Tr\Ham(s)\,ds$
  with $\widetilde\Ham(X)=\Ham(x)/\Tr\Ham(x)$ one has
  $\int \Tr\widetilde\Ham(X)\,dX = \int \Tr\Ham(x)\,dx$.  Thus finiteness
  or infiniteness of the total trace \emph{cannot} be altered by
  renormalising; in particular one cannot escape the tautology of item~(2)
  by ``choosing a different normalisation''.  The only genuine freedom is
  in specifying a \emph{different Hamiltonian} --- the putative normalised
  system $\Ham_\Phi$ appearing in the proposed Core Lemma
  \cref{eq:coreL}, with its own Weyl function --- rather than
  reparametrising a fixed one.
\end{remark}


%% ─────────────────────────────────────────────────────────────────────────────
\section{A trivial log-derivative realization}
\label{sec:trivial}

We first record that a canonical realization ``attached to $H_t$'' exists
by an entirely standard construction --- and precisely because it is so
cheap, it does not advance the programme.

\begin{proposition}[Log-derivative Herglotz function]\label{prop:trivial}
  Fix $t > \Lambda$.  Then $H_t$ is a real entire function in the
  Laguerre--P\'olya class, with Hadamard factorisation
  \begin{equation}\label{eq:hadamard}
    H_t(z) \;=\; C_t\, z^{2m}\, e^{-a_t z^2}
      \prod_{n}\Bigl(1 - \frac{z^2}{x_{n,t}^2}\Bigr)^{m_n},
    \qquad a_t \ge 0,\quad \sum_n \frac{m_n}{x_{n,t}^2} < \infty,
  \end{equation}
  where $x_{n,t}>0$ enumerate the \emph{positive} nonzero zeros of $H_t$
  with integer multiplicities $m_n\in\mathbb Z_{\ge1}$, and
  $m\in\mathbb Z_{\ge0}$ with $2m$ the order of the zero at the origin.
  Consequently the logarithmic derivative
  \begin{equation}\label{eq:mt}
    m_t(z) \;:=\; -\frac{H_t'(z)}{H_t(z)}
      \;=\; 2a_t z - \frac{2m}{z}
        + \sum_n m_n\Bigl(\frac{1}{x_{n,t}-z} + \frac{1}{-x_{n,t}-z}\Bigr),
  \end{equation}
  the series converging locally uniformly on
  $\mathbb C\setminus Z(H_t)$, is a Herglotz function:
  $\operatorname{Im}m_t(z) > 0$ for $\operatorname{Im}z > 0$.
\end{proposition}

\begin{proof}
  The factorisation \cref{eq:hadamard} is the standard form for the
  Laguerre--P\'olya class (a locally uniform limit of real polynomials
  with only real roots; see \cite{Newman1976} and, for the specific
  heat-flow zeros, \cite{KiKimLee2009,RodgersTao2020}).  Differentiating
  $\log H_t$ term by term (justified by the local uniform convergence
  from $\sum_n m_n/x_{n,t}^2 < \infty$) gives \cref{eq:mt}, each zero pair
  $\pm x_{n,t}$ contributing exactly once.  For $\operatorname{Im}z>0$
  each term $1/(x-z)$ with $x\in\R$ has
  $\operatorname{Im}\bigl(1/(x-z)\bigr) = \operatorname{Im}z/|x-z|^2 > 0$,
  the term $-2m/z$ contributes $2m\,\operatorname{Im}z/|z|^2 \ge 0$, and
  $2a_t z$ contributes $2a_t\operatorname{Im}z \ge 0$.  Since $H_t$ is
  nonconstant with at least one nonzero real zero, at least one summand
  contributes strictly, so $\operatorname{Im}m_t(z) > 0$ (not merely
  $\ge0$).
\end{proof}

\begin{corollary}[Automatic --- and vacuous --- limit-point]\label{cor:trivial}
  Applying Remling's inverse spectral theorem
  \cite[Theorem~5.1]{Remling2018} to $m_t$ of \cref{eq:mt} yields a
  \emph{unique trace-normalised} half-line canonical system
  $\widehat\Ham_t$ with $\Tr\widehat\Ham_t \equiv 1$ a.e.  Because $m_t$
  is a non-constant Herglotz function (indeed $\operatorname{Im}m_t>0$
  strictly, and the representing measure of $m_t$ is supported on
  infinitely many points), the associated system is not a single
  indivisible interval, so the standing assumptions of \cref{prop:weyl}
  hold.  Since
  $\Tr\widehat\Ham_t\equiv1$, $\int_0^\infty \Tr\widehat\Ham_t\,dx =
  +\infty$, and \cref{prop:weyl} places $x=+\infty$ in the limit-point
  case.

  Separately --- and this is a distinct observation --- setting
  $A_t = H_t$, $B_t = -H_t'$ gives $B_t/A_t = m_t$, so
  $E_t = A_t - iB_t = H_t + iH_t'$ is a Hermite--Biehler function.
  Indeed, since $H_t$ is real entire, $E_t^\# := \overline{H_t(\bar z)} -
  i\,\overline{H_t'(\bar z)} = H_t(z) - iH_t'(z)$, so
  $|E_t(z)|^2 - |E_t^\#(z)|^2 = 4|H_t(z)|^2\operatorname{Im}m_t(z) > 0$
  for $\operatorname{Im}z>0$ by \cref{prop:trivial}.  Zero-freeness of
  $E_t$ in the upper half-plane follows: if $E_t(z_0)=0$ and
  $H_t(z_0)\neq0$ then $\operatorname{Im}m_t(z_0) = -1 < 0$,
  contradicting the Herglotz property; if $H_t(z_0)=0$ then
  $\operatorname{Im}z_0 \le 0$ since all zeros of $H_t$ are real for
  $t>\Lambda$.  This does \emph{not} identify $E_t$ as a terminal
  $E$-function of the reconstructed half-line system $\widehat\Ham_t$; it
  is only a companion observation about $H_t$.
\end{corollary}

\begin{remark}[Why this does not help]\label{rem:trivial-weak}
  The limit-point conclusion of \cref{cor:trivial} is a tautological
  consequence of trace normalisation, not a fact about the geometry of a
  ``natural'' Hamiltonian: by \cref{rem:invariance} every
  trace-normalised half-line system has infinite total trace, so
  \emph{every} such realization is limit-point regardless of $H_t$.  The
  construction uses only that $H_t$ has real zeros with the summability
  \cref{eq:hadamard}; it sees nothing arithmetic and carries no
  information toward RH.  In particular it does \emph{not} exhibit $H_t$
  as the terminal $A$-function of a prescribed canonical chain.
\end{remark}

%% ─────────────────────────────────────────────────────────────────────────────
\section{Audit of the programme's Core Lemma}
\label{sec:open}

The trivial construction of \cref{sec:trivial} shows that producing
\emph{some} canonical system with $H_t$ in view is automatic.  The
programme \cite{SpecRoadmap} instead proposes a specific spectral
criterion, its ``Core Lemma~L'', which we now audit.  We find its stated
justification \emph{insufficient as written}, and we record a corrected
--- necessarily weaker, one-directional --- replacement.

\subsection{The proposed criterion, and why its justification is insufficient}

The roadmap \cite{SpecRoadmap} states its D1 acceptance criterion as a two-part
equivalence.  In our notation:
\begin{equation}\label{eq:coreL}
  \begin{aligned}
  \text{RH}\;\stackrel{?}{\Longleftrightarrow}\;
  &\bigl(T_\Phi\text{ is limit-point at }+\infty\bigr)\\
  &\quad\text{and}\quad
  \bigl(\text{a fixed self-adjoint realization has }
        \operatorname{spec}=\{\gamma_n\}\bigr),
  \end{aligned}
\end{equation}
where $T_\Phi$ is a symmetric relation on a half-line canonical system
built (in \cite{SpecRoadmap}) from the fixed positive kernel $\Phi$
through the cosine-transform relation F1, and $\{\gamma_n\}$ are the
imaginary parts of the nontrivial zeros of $\zeta$.  We write
$\stackrel{?}{\Longleftrightarrow}$, not $\iff$, deliberately.

\begin{remark}[The proposed implication is not justified as
written]\label{rem:coreL-defective}
  A self-adjoint relation has \emph{real} spectrum.  If some nontrivial
  zero were off the critical line, $\rho_0=\beta+i\gamma$ with
  $\beta\neq\tfrac12$, its ordinate $\gamma$ would still be a real number.
  A statement about the real spectrum $\{\gamma_n\}$ therefore records the
  ordinates only and is \emph{blind to the real parts} $\beta$.  Hence
  ``$\operatorname{spec}=\{\gamma_n\}$'', \emph{by itself}, does not imply
  $\beta=\tfrac12$ and does not imply RH.

  We are careful about the scope of this observation.  It shows that
  self-adjointness together with equality of ordinate sets is
  \emph{insufficient} to force RH, and hence that the roadmap's stated
  justification is incomplete.  It does \emph{not} disprove the
  right-to-left implication for the particular relation $T_\Phi$: since
  $T_\Phi$ has not been rigorously defined (\cref{rem:undefined}), it
  remains logically possible that a specific construction from the full
  $\Phi/\xi$ data would encode the real parts indirectly, e.g.\ via a
  theorem
  $\operatorname{spec}T_\Phi=\{\gamma_n\}\Rightarrow
   \operatorname{div}D_{\theta_0}=\operatorname{div}\Xi$.  What is
  missing is exactly such a theorem: an additional argument recovering the
  real-part information, such as the divisor identity \cref{eq:divid}.
  Accordingly \cref{eq:coreL} as written is not a justified restatement of
  RH, and the Core Lemma needs this gap filled before it can be attacked.
\end{remark}

\subsection{The corrected statement: a one-way sufficiency}

The object that \emph{does} retain the real parts is the entire function
\begin{equation}\label{eq:Xi}
  \Xi(z) \;=\; \xi\!\left(\tfrac12+iz\right).
\end{equation}
A nontrivial zero $\rho=\beta+i\gamma$ of $\zeta$ corresponds to the zero
$z=\gamma+i(\tfrac12-\beta)$ of $\Xi$, which is \emph{real} iff
$\beta=\tfrac12$.  Here $\operatorname{div}$ denotes the zero divisor,
i.e.\ the formal sum of zeros counted with their orders.  Thus
real-rootedness of $\Xi$ is equivalent to RH, and a spectral model implies
RH only if it pins down the \emph{full zero divisor} of $\Xi$, not merely
a set of real ordinates.

\begin{problem}[Repair obligations for D1]\label{prob:gen}
  To turn the programme's Core Lemma into a well-posed, one-directional
  sufficiency for RH, the following four obligations must be discharged;
  each has an explicit PASS condition.
  \begin{enumerate}[label=(\arabic*)]
    \item \textbf{Construction (non-circular).}  Specify a mathematically
          checkable map $\Phi\mapsto\Ham_\Phi$, including the exact F1
          identity and the normalisation, yielding a locally integrable,
          positive-semidefinite
          $\Ham_\Phi\colon[0,\infty)\to M_2(\R)$ and the minimal symmetric
          relation $T_\Phi$ it induces via \cref{eq:cansys}.  \emph{PASS
          condition.}  The
          datum F1 and the rule $\Phi\mapsto\Ham_\Phi$ must be fixed
          \emph{independently} of the zero divisor of $\Xi$, of the
          ordinates $\{\gamma_n\}$, and of RH: neither the desired
          spectrum nor any measure manufactured from it may enter as
          input.  Existence \emph{and uniqueness} of the resulting
          normalised $\Ham_\Phi$ must be proved --- if one $\Phi$ could
          yield different $\Ham_\Phi$ under different admissible choices,
          ``the specific operator $T_\Phi$'' would still be undefined.
    \item \textbf{Standing assumptions and length.}  Prove that
          $\Ham_\Phi$ meets the standing assumptions of \cref{prop:weyl}
          and, \emph{independently}, that
          $\int_0^\infty\Tr\Ham_\Phi(x)\,dx=+\infty$ (whence limit-point,
          by \cref{prop:weyl}).
    \item \textbf{Discreteness and characteristic function
          (constructed, not named).}  After fixing $\theta_0\in[0,\pi)$,
          prove that the resulting self-adjoint realization has
          \emph{purely discrete} spectrum --- with no absolutely
          continuous or singular continuous part and no finite real
          accumulation point.  Then construct $D_{\theta_0}$ by a
          \emph{prescribed operator-theoretic procedure} from $T_\Phi$ and
          the boundary condition --- a boundary/Wronskian determinant, a
          regularised Fredholm determinant of the resolvent, or a
          rigorously convergent half-line limit of truncated transfer
          determinants --- and prove
          \begin{equation}\label{eq:charzero}
            D_{\theta_0}(\lambda)=0
            \iff \lambda\in\sigma_{\mathrm p}(T_{\Phi,\theta_0}),
          \end{equation}
          with the order of each zero equal to the (separately specified)
          spectral multiplicity, and with the real normalisation
          $D_{\theta_0}^{\#}=D_{\theta_0}$.  \emph{PASS condition.}
          $D_{\theta_0}$ may \emph{not} be defined as a canonical product
          over prescribed points (in particular not as $\Xi$, nor over
          $\{\gamma_n\}$); it must arise from $T_\Phi$ by the fixed
          procedure above.
    \item \textbf{Divisor identity.}  Prove the zero-divisor identity
          \begin{equation}\label{eq:divid}
            \operatorname{div} D_{\theta_0}
            \;=\;
            \operatorname{div}\Xi .
          \end{equation}
  \end{enumerate}
  Given (1)--(4): since the eigenvalues of a self-adjoint realization are
  real, \cref{eq:divid} forces all zeros of $\Xi$ to be real, i.e.\ RH.
  The point of the PASS conditions in (1) and (3) is exactly that only
  then is \cref{eq:divid} a genuine \emph{theorem} rather than a
  definition of the objects in terms of the target.
\end{problem}

\begin{remark}[Direction, and difficulty, of \cref{prob:gen}]\label{rem:coreL-is-RH}
  Only the implication ``(1)--(4)$\Rightarrow$RH'' is claimed; the
  converse is not, and RH does not by itself force any particular
  $T_\Phi$ to have characteristic function $\Xi$.  Thus \cref{prob:gen} is
  \emph{at least as hard as} RH, and it imposes additional requirements
  not known to follow from RH: not only that the zeros of $\Xi$ be real
  but that they coincide, with order, with the eigenvalue divisor of a
  \emph{specific} operator.
  Contrary to a claim in an earlier draft of this audit, this does not
  mean there is no RH-independent acceptance criterion: the identity
  \cref{eq:divid} is perfectly well-posed to \emph{state} without knowing
  RH, and its target $\operatorname{div}\Xi$ is explicit.  What is missing
  is different --- a rigorous definition of $T_\Phi$ and $D_{\theta_0}$,
  and a proof relating the latter to the already explicit target divisor
  $\operatorname{div}\Xi$ (\cref{rem:undefined}).  This paper resolves
  none of it.
\end{remark}

\begin{remark}[Three distinct spectral conditions, not one]\label{rem:three-conditions}
  ``$\operatorname{spec}=\{\gamma_n\}$ with multiplicities and correct
  weights'', as an earlier draft put it, conflates three genuinely
  different requirements that must be separated:
  \begin{enumerate}[label=(\alph*)]
    \item \emph{support}: $\operatorname{supp}\rho=\{\gamma_n\}$, using the
          cyclic Weyl spectral measure $\rho$ and its closed support (a
          statement about the spectrum as a set);
    \item \emph{spectral multiplicity}: the eigenvalue multiset, i.e.\
          the multiplicity of each point of the spectrum.  Under
          additional hypotheses ensuring simple eigenvalues (which a
          separated boundary condition on a scalar half-line system
          typically, but not automatically, provides), matching
          zeta-zero orders would further force the zeros of $\Xi$ to be
          \emph{simple} --- a statement strictly stronger than RH;
    \item \emph{norming constants}: the weights $\rho(\{\gamma_n\})$ of
          the atoms of the spectral measure, which are \emph{not}
          eigenvalue multiplicities but the squared reciprocal norms of
          the eigenfunctions, after fixing the normalisation of the
          canonical eigenfunctions.
  \end{enumerate}
  None of (a)--(c) is equivalent to another.  The support condition (a)
  alone is insufficient; the divisor identity \cref{eq:divid} used here
  does retain the missing real-part information (other spectral data
  encoding the real parts could serve equally).  No formula for the
  weights in (c) is available in \cite{SpecRoadmap}; that gap is part of
  \cref{rem:undefined}.
\end{remark}

\begin{remark}[$T_\Phi$ and $\Ham_\Phi$ are not yet defined]\label{rem:undefined}
  The roadmap \cite{SpecRoadmap} refers to $T_\Phi$ ``induced by $\Phi$ through F1''
  but supplies no Hamiltonian formula, transfer matrix, Weyl function, or
  de~Branges chain.  Consequently one cannot yet speak rigorously of the
  minimal/maximal relation of $T_\Phi$, its self-adjoint extensions, or
  its spectral measure: these presuppose $\Ham_\Phi$.  The genuinely first
  task --- logically prior to any spectral identification --- is therefore
  the \emph{construction} of $\Ham_\Phi$ and $T_\Phi$, stated as
  obligation~(1) of \cref{prob:gen}.  The target $\operatorname{div}\Xi$
  is, by contrast, already explicit; what is missing is a rigorous
  definition of $T_\Phi$ and $D_{\theta_0}$ and a proof relating the
  latter to that target.  Until obligation~(1) is done, \cref{eq:coreL}
  and the left-hand side of \cref{eq:divid} refer to objects that are not
  defined, and the phrases
  ``the Weyl function of $T_\Phi$'' or ``the spectral measure $\rho$'' are
  placeholders.
\end{remark}

\begin{remark}[Why the log-derivative realization does not solve this]\label{rem:not-trivial}
  \Cref{cor:trivial} produces \emph{a} half-line system whose spectral
  measure is atomic, supported on the zeros of $H_t$.  Even setting aside
  that it works from $H_t$ ($t>\Lambda$) rather than the putative system
  constructed from the undeformed kernel $\Phi$, it does not solve
  \cref{prob:gen}: it matches a \emph{zero set} by construction, whereas
  \cref{prob:gen} demands the zero-divisor identity \cref{eq:divid} for a
  \emph{specific}, independently constructed $T_\Phi$ (and, by the PASS
  condition of \cref{prob:gen}(3), for a $D_{\theta_0}$ produced from
  $T_\Phi$, not defined over prescribed points).  Matching a prescribed
  set of points by choosing a Herglotz function is trivial; proving that a
  \emph{given} operator has a prescribed characteristic divisor is the
  whole problem.
\end{remark}

\begin{remark}[Why no terminal entire function is imposed]\label{rem:incompat}
  The D1 model of \cite{SpecRoadmap} is a singular canonical system on
  $[0,\infty)$.  The functions $E_L$ associated with finite truncations
  $[0,L]$ are approximation data for this half-line system; D1 does not
  assume that $E_L$ converges to a nonzero entire function as
  $L\to\infty$.  In particular the condition $A_\infty = c\,\Xi$ is
  neither required nor used: a nonzero entire terminal function is
  characteristic of a \emph{finite, regular} chain, whereas a genuine
  half-line limit-point endpoint typically admits no such terminal entire
  function.  (The characteristic function $D_{\theta_0}$ of
  \cref{eq:divid}, when it exists, is a different object --- attached to
  the self-adjoint realization, not a terminal transfer-matrix entry.)
  Imposing $A_\infty=c\,\Xi$ would convert D1 into a separate finite-chain
  realization problem, which the programme does not require.
\end{remark}

\begin{remark}[The limit-point conjunct is conditional, not settled]\label{rem:length}
  The first conjunct of \cref{eq:coreL} --- limit-point at $+\infty$ ---
  is often described in \cite{SpecRoadmap} as the ``easy half'', on the
  ground that the Riemann--von~Mangoldt law
  $N(T)=\tfrac{T}{2\pi}\log\tfrac{T}{2\pi e}+O(\log T)$ ``forces infinite
  canonical length''.  We do \emph{not} endorse that inference.  $N(T)$
  counts zeros in the \emph{spectral} variable, whereas
  $\int_0^\infty\Tr\Ham_\Phi\,dx$ is a length in the \emph{spatial}
  variable.  Such a deduction would require \emph{either} a direct
  estimate of the total trace from an explicit formula for $\Ham_\Phi$
  (needing no spectral identification at all), \emph{or} a proved
  spectral-counting correspondence --- weaker than the full divisor
  identity \cref{eq:divid} --- combined with a suitable
  Cartwright/de~Branges growth/type theorem converting spectral density
  into canonical length.  (We do not assert that
  $N(T)\sim\tfrac{T}{2\pi}\log T$ alone yields trace divergence; a precise
  such theorem would have to be cited and its hypotheses checked.)
  Neither is supplied here.  Moreover infinite
  \emph{discrete} spectrum alone does not force infinite length --- a
  finite regular canonical system can have infinite discrete spectrum ---
  so the super-linear $T\log T$ growth would have to be used explicitly to
  exclude finite exponential type.  Accordingly we treat trace divergence
  as a \emph{separate hypothesis}, not a settled fact:
  \begin{quote}
    No theorem translating the Riemann--von~Mangoldt counting law into
    $\int_0^\infty\Tr\Ham_\Phi(x)\,dx=\infty$ is proved in this paper.
  \end{quote}
  Even granting it, by \cref{rem:invariance} it is a reparametrisation
  invariant that cannot be manufactured by renormalisation, and it must
  not be re-derived from \cref{prop:weyl} on pain of circularity.
\end{remark}

\scopenote{\textbf{OPEN.}  The following obstacles constrain any
resolution of \cref{prob:gen}; none is resolved here.  They also explain
why the heat-flow family $H_t$ ($t>\Lambda$) of \cref{sec:trivial} ---
\cite{SpecRoadmap}'s $\Lambda$-1/$\Lambda$-2 step, where all zeros are
already real --- gives no RH-independent handle on the putative system
constructed from the undeformed kernel $\Phi$.}

\medskip\noindent\textbf{Obstacle~1: a real entire function is never
Hermite--Biehler.}
For any real entire $F$ one has $F^\# = F$, so the strict de~Branges
inequality $|E(z)|>|E^\#(z)|$ ($\operatorname{Im}z>0$) fails identically.
Thus neither $\Xi$ nor $H_t$ can itself be an $E$-function; the spectral
data of the half-line system must be encoded through a genuine
Hermite--Biehler $E=A-iB$ (equivalently a Weyl function $m$), not by
declaring a real entire function a structure function.

\medskip\noindent\textbf{Obstacle~2: the naive $|H_t|^{-2}\,d\tau$ is not
a spectral measure.}
For $t>\Lambda$ all zeros of $H_t$ are real, and there are infinitely
many \cite{KiKimLee2009,RodgersTao2020}.  At each real zero
$|H_t(\tau)|^{-2}$ has a non-integrable double pole, so
$d\mu_t = |H_t(\tau)|^{-2}\,d\tau$ is not locally finite; it cannot be the
representing measure of any Weyl function.

\medskip\noindent\textbf{Obstacle~3: no controlled spectral-system limit
toward the candidate parameter $t=0$.}
Even a complete identification for $H_t$ at a fixed $t>\Lambda$ would
carry no RH information: there the zeros of $H_t$ are already all real
(\cite{SpecRoadmap} $\Lambda$-2), so real-rootedness is not in question.
The programme locates RH at the limit down to the candidate parameter
$t=0$ (the true critical parameter is $\Lambda$, not known to be $0$).
Here $t$ is the heat-flow parameter of \cref{eq:heatflow}; the roadmap's
symbol $\lambda$ denotes its flow parameter, which agrees with $t$ up to
the fixed rescaling recorded after \cref{eq:heatflow}.  We use $t$
throughout.  The roadmap's constraint $(C')$ argues that the relevant
minimum should tend to $0$ as $t\downarrow0^+$ with \emph{equality}: a
uniform positive margin would, \emph{provided} the requisite continuity
of that minimum in the parameter, its strict equivalence with the
real-zero property, and a stability/continuation theorem past $0$ are
established, extend the conclusion below $0$ and hence contradict
$\Lambda\ge0$.  Those implications are the roadmap's, not results proved
here.  Note that the function $H_t$ of course has a limit as
$t\downarrow0$; what is missing is any control of the \emph{canonical
system or spectral data} along that limit.  No such control is provided
here.

%% ─────────────────────────────────────────────────────────────────────────────
\section{The limit-point criterion (reference statement)}
\label{sec:L3}

\begin{proposition}[Trace-divergence limit-point criterion]\label{prop:weyl}
  Let $\Ham \colon [0,\infty) \to M_2(\R)$ satisfy the standing
  assumptions of \cite{Remling2018}: $\Ham \in L^1_{\mathrm{loc}}$,
  $\Ham \succeq 0$, $\Ham(x)\neq 0$ for a.e.\ $x$, and $[0,\infty)$ is not
  a single indivisible interval.  Then the right endpoint $x=+\infty$ of
  the canonical system \cref{eq:cansys} is in the limit-point case if and
  only if
  \[
    \int_0^{\infty} \Tr\Ham(x)\,dx \;=\; +\infty.
  \]
  This is a specialisation of \cite[Theorem~3.5]{Remling2018}; it is
  stated here only for reference.
\end{proposition}

\begin{remark}[On the standing assumptions]\label{rem:definite}
  The a.e.\ non-vanishing and no-single-indivisible-interval hypotheses
  are Remling's standing conventions.  They do \emph{not} by themselves
  turn the maximal/minimal objects into operators: even with
  $\Ham(x)\neq0$ a.e., rank-one (indivisible) subintervals force one to
  work with linear \emph{relations}, which is why we use relation language
  throughout (cf.\ \cite[Theorem~2.25]{Remling2018}).  A merely
  positive-semidefinite $\Ham$ vanishing on a positive-measure set is
  reduced to this setting by the trace reparametrisation of
  \cref{rem:invariance}; by that same invariance the trace integral, and
  hence the classification, is unchanged.
\end{remark}

\begin{proof}
  By \cite[Theorem~3.5]{Remling2018}, under the standing assumptions the
  right endpoint is limit-point iff
  $\Ham \notin L^1((c,\infty);M_2)$ for one --- equivalently every ---
  $c>0$.  Since $\Ham \succeq 0$, positivity gives $\Ham_{11},\Ham_{22}\ge0$
  and $|\Ham_{12}| \le \sqrt{\Ham_{11}\Ham_{22}} \le \tfrac12\Tr\Ham$,
  whence for each $c>0$
  \[
    \Ham \in L^1((c,\infty);M_2)
    \iff
    \Tr\Ham \in L^1(c,\infty).
  \]
  Combining the two equivalences gives the stated trace form.  (The
  reference uses $J y' = -z\Ham y$; the substitution $z\mapsto -z$ gives
  the sign convention of \cref{eq:cansys} and does not affect endpoint
  classification.  The same criterion appears in Winkler
  \cite{Winkler1995}.)

  For the reader's convenience we indicate the direction that has an
  elementary proof.  If $\int_c^\infty \Tr\Ham < \infty$ then
  $\Ham \in L^1(c,\infty)$; writing \cref{eq:cansys} as the integral
  equation $y(x) = y(c) + z\!\int_c^x J^{-1}\Ham(s)y(s)\,ds$ and applying
  Gronwall's inequality, every solution at every fixed $z$ is bounded
  near $+\infty$, so
  $\int_c^\infty y(x,z)^*\Ham(x)y(x,z)\,dx
   \le (\sup\|y\|^2)\int_c^\infty \Tr\Ham < \infty$;
  all solutions are then in $L^2_\Ham$ and $+\infty$ is limit-circle.
  The reverse implication --- trace divergence $\Rightarrow$ limit-point
  --- is \emph{not} elementary in this way: exhibiting a single non-$L^2$
  solution at the special value $z=0$ (the constant vector spanning the
  larger of $\int_c^\infty\Ham_{jj}$) does not by itself control the
  solution count at non-real $z$.  That step requires the parameter
  independence of the deficiency index, which is precisely the nontrivial
  content of \cite[Theorem~3.5]{Remling2018} via
  \cite[Lemma~3.7]{Remling2018}; we do not reproduce it, and rely on the
  cited theorem for the full equivalence.
\end{proof}

%% ─────────────────────────────────────────────────────────────────────────────
\section{Conditional corollary and boundary structure}
\label{sec:proof}

\begin{corollary}[The limit-point conjunct, conditionally]\label{cor:conditional}
  Assume the construction of \cref{rem:undefined} is carried out, giving a
  half-line Hamiltonian $\Ham_\Phi$ and relation $T_\Phi$ satisfying the
  standing assumptions of \cref{prop:weyl}, and assume the separate
  hypothesis $\int_0^\infty \Tr\Ham_\Phi(x)\,dx = +\infty$ of
  \cref{rem:length}.  Then $x=+\infty$ is in the limit-point case.  This is
  the first conjunct of \cref{eq:coreL}; it is \emph{conditional} on both
  the construction and the trace-divergence hypothesis, the latter to be
  supplied from the construction and \emph{not} from \cref{prop:weyl}
  (\cref{rem:length}).
\end{corollary}

\begin{proof}
  Immediate from \cref{prop:weyl} applied to $\Ham_\Phi$.
\end{proof}

\begin{remark}[This is not unconditional, and not the trivial case]\label{rem:not-uncond}
  \Cref{cor:conditional} settles only the first conjunct of
  \cref{eq:coreL}, and only modulo the construction and the
  trace-divergence input.  The second conjunct is the corrected
  identification \cref{prob:gen}: even in its right form (the divisor
  identity \cref{eq:divid}) it is at least as hard as RH
  (\cref{rem:coreL-is-RH}), and it is \emph{not} addressed.  For the
  trace-normalised realization of \cref{cor:trivial} the trace-divergence
  hypothesis is automatic but yields only the tautology of
  \cref{rem:trivial-weak}.  So the paper establishes \emph{nothing}
  arithmetically meaningful: all content sits in the spectral
  identification, which we do not resolve.
\end{remark}

\begin{corollary}[Boundary-parameter structure, under
\cref{cor:conditional}]\label{cor:alpha}
  Under the hypotheses of \cref{cor:conditional}, the minimal symmetric
  \emph{relation} $T_\Phi$ associated with $\Ham_\Phi$ has deficiency
  indices $(1,1)$: the right endpoint $x=+\infty$ is limit-point and
  contributes no free parameter, while the regular left endpoint $x=0$
  carries one self-adjoint boundary condition.  For every
  $\theta_0\in[0,\pi)$ the condition $e_{\theta_0}^{*}J\,y(0)=0$, with
  $e_{\theta_0} = (\cos\theta_0,\sin\theta_0)^{\top}$, defines a
  self-adjoint realization, and these exhaust all self-adjoint
  realizations \cite[Theorem~2.25]{Remling2018}.  In particular
  limit-point at $+\infty$ does \emph{not} make $T_\Phi$ essentially
  self-adjoint and does \emph{not} yield a \emph{unique} self-adjoint
  realization; this corrects the ``$\alpha$ does not appear / unique
  self-adjoint realization'' phrasing of \cite{SpecRoadmap}'s D1, which
  overstates the consequence of limit-point.

  We phrase this for relations throughout: because indivisible
  subintervals may occur even when $\Ham_\Phi(x)\neq0$ a.e.\ (cf.\
  \cref{rem:definite}), the maximal/minimal objects need not be operators.
  Upgrading to a densely defined operator requires separately showing the
  multivalued part of the maximal relation is trivial; this is governed by
  the operator-part / multivalued-part decomposition of the relation (see
  the corresponding discussion in \cite{Remling2018}), not by
  \cite[Theorem~2.25]{Remling2018}, which concerns only the parametrisation
  of self-adjoint boundary conditions.
\end{corollary}

%% ─────────────────────────────────────────────────────────────────────────────
\section{Discussion and scope}
\label{sec:discussion}

\subsection{What this paper establishes}

The only proved statements are (i)~the reference criterion
\cref{prop:weyl}, a direct specialisation of
\cite[Theorem~3.5]{Remling2018}, (ii)~the trivial log-derivative
realization \cref{prop:trivial}--\cref{cor:trivial}, and (iii)~the
\emph{negative} observation that the programme's Core Lemma
\cref{eq:coreL} is \emph{not justified as written} as an equivalence with
RH (\cref{rem:coreL-defective}).  None is a new theorem in analysis:
(i)~is standard and depends on \cite{Remling2018}; (ii)~is a routine
consequence of the Laguerre--P\'olya factorisation; (iii)~is an
elementary but consequential identification of a logical gap.  We
therefore present the paper as a reference note plus a logical-gap audit,
not a source of new theorems.

We do \emph{not} establish essential self-adjointness of any operator (or
relation) attached to the de Bruijn--Newman flow, nor a unique
self-adjoint realization, nor any arithmetically meaningful limit-point
property of that flow, for any value of $t$.

\subsection{What remains open}

The corrected identification problem \cref{prob:gen} is open.  By
\cref{rem:coreL-is-RH} only the one-way sufficiency
``$\operatorname{div}D_{\theta_0}=\operatorname{div}\Xi\Rightarrow$RH''
holds, and it is at least as hard as RH --- imposing additional
requirements not known to follow from RH, since it also asks that the
divisor match that of a \emph{specific} operator
(\cref{rem:three-conditions}).  Logically prior to it is the
unstated \emph{construction} of $\Ham_\Phi$ and $T_\Phi$
(\cref{rem:undefined}); the target $\operatorname{div}\Xi$ is by contrast
already explicit.  Until the construction is done, $T_\Phi$,
$D_{\theta_0}$, the Weyl function, and the spectral measure --- i.e.\ the
left-hand data of \cref{eq:coreL,eq:divid} --- refer to objects that are
not yet defined.

The limit-point conjunct of \cref{eq:coreL} is \emph{not} the hard part,
but neither is it settled: as recorded in \cref{rem:length}, no theorem
translating the Riemann--von~Mangoldt law into
$\int_0^\infty\Tr\Ham_\Phi\,dx=\infty$ is proved here, so trace divergence
remains a hypothesis --- which, if granted, gives limit-point via
\cref{prop:weyl}, subject to the anti-circularity and invariance cautions
of \cref{rem:invariance,rem:length}.  We also correct
\cite{SpecRoadmap}'s D1 write-up on one point of principle: limit-point at
$+\infty$ does not give a \emph{unique} self-adjoint realization --- the
regular left endpoint retains the parameter $\theta_0$ (\cref{cor:alpha}).

Finally, even a correct divisor identity \cref{eq:divid} presupposes that
the realization has discrete spectrum and a characteristic entire
function; a general half-line self-adjoint canonical system need have
neither (it may have absolutely continuous or singular continuous
spectrum), so establishing purely discrete spectrum with no continuous
part is itself part of the open problem.  None of this is addressed here,
and the critical limit (\cref{rem:length}, Obstacle~3 of \cref{sec:open})
lies outside the present scope.

\subsection{Aside: the Conrey--Li obstruction and finite-chain
realizations}

This aside is \emph{not} part of the D1 half-line problem; it concerns the
separate finite-chain realization that D1 does not require.  The
Conrey--Li work \cite{ConreyLi2000} concerns specific positivity
conditions of de~Branges type for $\xi$/$L$-functions; it is \emph{not}
merely the observation that the real entire candidate
$E_T(z)=\xi(1/2-iz)$ fails the Hermite--Biehler test (that failure is
trivial, since $E_T^\#=E_T$).  A finite-chain de~Branges realization with
a terminal $A$-function proportional to $H_t$ would require building a
genuine $E_t=A_t-iB_t$ with nonzero $B_t$; this is a separate problem and
is \emph{not} needed for D1 (\cref{rem:incompat}).  The precise relation
of any such $E_t$ to the Conrey--Li positivity conditions --- for which we
refer the reader to \cite{ConreyLi2000} directly --- is beyond this audit.
Any earlier claim that $E_\lambda = \Xi_\lambda$ ``avoids'' the
obstruction was mistaken.

%% ─────────────────────────────────────────────────────────────────────────────
\begin{thebibliography}{99}

\bibitem{deBranges1968}
L.~de~Branges,
\emph{Hilbert Spaces of Entire Functions},
Prentice-Hall, Englewood Cliffs, NJ, 1968.

\bibitem{Krein1952}
M.~G.~Kre\u{\i}n,
On the indeterminate case of the Sturm--Liouville boundary problem in the
interval $(0,\infty)$,
\emph{Izv.\ Akad.\ Nauk SSSR Ser.\ Mat.} \textbf{16} (1952), 293--324.

\bibitem{Remling2018}
C.~Remling,
\emph{Spectral Theory of Canonical Systems},
De Gruyter Studies in Mathematics~70, Berlin, 2018.

\bibitem{Winkler1995}
H.~Winkler,
The inverse spectral problem for canonical systems,
\emph{Integral Equations Operator Theory} \textbf{22} (1995), 360--374.

\bibitem{ConreyLi2000}
J.~B.~Conrey and X.-J.~Li,
A note on some positivity conditions related to zeta and $L$-functions,
\emph{Int.\ Math.\ Res.\ Not.} \textbf{2000}, no.~18, 929--940.

\bibitem{RodgersTao2020}
B.~Rodgers and T.~Tao,
The de Bruijn--Newman constant is non-negative,
\emph{Forum Math.\ Pi} \textbf{8} (2020), e6.

\bibitem{Newman1976}
C.~M.~Newman,
Fourier transforms with only real zeros,
\emph{Proc.\ Amer.\ Math.\ Soc.} \textbf{61} (1976), 245--251.

\bibitem{KiKimLee2009}
H.~Ki, Y.-O.~Kim, and J.~Lee,
On the de Bruijn--Newman constant,
\emph{Adv.\ Math.} \textbf{222} (2009), 281--306.

\bibitem{SaouterGourdon2011}
Y.~Saouter, X.~Gourdon, and P.~Demichel,
An improved lower bound for the de Bruijn--Newman constant,
\emph{Math.\ Comp.} \textbf{80} (2011), 2281--2287.

\bibitem{PlattTrudgian2021}
D.~J.~Platt and T.~S.~Trudgian,
The Riemann hypothesis is true up to $3\cdot10^{12}$,
\emph{Bull.\ London Math.\ Soc.} \textbf{53} (2021), 792--797.

\bibitem{Titchmarsh1986}
E.~C.~Titchmarsh,
\emph{The Theory of the Riemann Zeta-Function},
2nd ed., revised by D.~R.~Heath-Brown, Oxford, 1986.

\bibitem{SpecRoadmap}
L.~Tao,
\emph{Speculative Roadmap: Critical Canonical System / $\Lambda{=}0$
Analysis}, internal document,
\texttt{docs/SPECULATIVE\_ROADMAP.md},
\texttt{weil\allowbreak-first\allowbreak-prime} repository, 2026.

\end{thebibliography}

\end{document}
